\section{Model \& Open Weight: ``公開''の先で何が違うのか}
\sectionkicker{Qwen3.8-Max / DeepSeek-V4-Flash-0731 / Kimi K3}

\textbf{W32で目立った3モデルは、単純な「新モデル3件」ではない。}
Qwen3.8-Max-Previewは7月19日に先行提供され、W32ではagent/coding評価とopen-weightへの期待が再浮上した。DeepSeek-V4-Flash-0731は7月31日のre-post-trainingによるAPI更新である。Kimi K3は巨大なfull weightsが既に公開されており、W32では「どこまで小さな実行環境へ持ち込めるか」という極端なcommunity experimentが注目された\cite{qwen38-token-plan,deepseek-v4-flash-0731,kimi-k3-hf,kimi-k3-x-lowram}。3件を並べると、frontier modelを追う軸がrelease日やbenchmark scoreだけでは足りなくなっていることが分かる。

\begin{tcolorbox}[enhanced,breakable,colback=SurveySoft,colframe=SurveyRule,boxrule=0.4pt,arc=1mm,left=2mm,right=2mm,top=1.5mm,bottom=1.5mm]
\textbf{3モデルのW32における位置づけ}\par\smallskip
\textbf{Qwen3.8-Max}: 7/19 Previewの再評価。W32ではagent/coding評価とopen-weightへの期待が前面に出た。\par
\textbf{DeepSeek-V4-Flash-0731}: 7/31のre-post-training/API更新。architecture・sizeはPreviewから据え置き。\par
\textbf{Kimi K3}: full weights公開後の運用段階。W32では巨大MoEをlocalでどう動かすかというsystems trade-offが注目された。\par\smallskip
\footnotesize これは性能順位ではなく、「何が今週技術的に重要になったか」の比較である。
\end{tcolorbox}

\subsection*{Qwen3.8-Max --- releaseではなく「次に何が開くか」への期待}

Alibaba Cloud Model Studioによれば、Qwen3.8-Max-Previewのdebutは\textbf{7月19日}であり、W32の新規releaseではない。Token Plan for Individualで先行提供され、full-stack development、data analysis、vision understandingなどを用途として挙げていた\cite{qwen38-token-plan}。

それでもW32で再び候補になったのは、X上でagent/coding評価やmultimodal testが共有され、さらにopen-weight化への期待が強く語られたためである。ここでは重要な境界がある。\textbf{W32で確認できたのはPreviewへの技術的関心であって、weightsやlicenseの正式公開ではない。}したがってQwenを「今週open-weightになったモデル」として扱うことはできない。

\begin{communitynote}[Qwen: benchmarkよりも「次に何が公開されるか」]
Reaction PassではCutoff前に3件の代表投稿が得られ、agent/coding評価、multimodal testing、open-weightへの期待が中心だった。公式Qwen投稿も含まれるため、rankや能力主張はvendor/social claimとして読む必要がある\cite{qwen38-x-ranking}。一部の実務移行を強く主張する投稿はCutoff後であり、本稿ではMain-window signalと分離する。
\end{communitynote}

\subsection*{DeepSeek --- architectureを変えず、post-trainingでagent側を更新}

DeepSeekのchange logは、7月31日に\texttt{deepseek-v4-flash} APIのpublic beta updateを記録している。DeepSeek自身の説明では、\textbf{DeepSeek-V4-Flash-0731はPreviewと同一architecture・同一sizeで、変更はre-post-trainingのみ}である。また今回の更新対象はFlash APIであり、V4-Pro APIやAPP/WEB modelは変更されていない\cite{deepseek-v4-flash-0731}。

これは「新しいbase modelを出した」というより、同じmodel bodyに対してagent capabilityを強く再post-trainしたeventとして読む方が正確である。DeepSeekはTerminal Bench 2.1、DeepSWE、Toolathlon等の数値を公開しているが、それらは最大effort、temperature 1.0、DeepSeek Harness minimal modeなどの条件を伴うvendor evaluationである\cite{deepseek-v4-flash-0731}。

公式model cardで示されるV4-Flash系は284B total / 13B activated、1M contextのMoEで、公開repositoryはMIT Licenseで提供されている\cite{deepseek-v4-flash-hf}。一方、今回Reaction Passで集まった0731版への代表的反応は\textbf{全件Cutoff後}だったため、W32 Main Community Watchの強い根拠には使わない。

\begin{claimboundary}[同じbenchmark名でもharnessが違えば単純比較できない]
DeepSeekのagent benchmark、Qwenに対するX上のranking、Kimi K3のcoding benchmarkは、modelだけでなくharness、reasoning effort、tool setup、sampling条件が異なる。したがって本稿では3モデルのscoreを一つの順位表へ統合しない。
\end{claimboundary}

\subsection*{Kimi K3 --- 「巨大weightsを公開する」と「手元で実用的に動く」は別問題}

Moonshot AIのofficial model cardはKimi K3をopen-weight native multimodal agentic modelと位置づける。総parameter数は2.8T、activated parameterは104B、context lengthは1,048,576 tokensで、KDAとAttention Residualsを採用する。weightsはMXFP4、activationはMXFP8を用いるquantization-aware training構成で、full weightsはKimi K3 Licenseの下で公開されている\cite{kimi-k3-hf}。

ここで「open weights」の意味を現実の運用へ戻すと、Hugging Face repository自体がおよそ\textbf{1.56 TB}ある\cite{kimi-k3-files}。つまり公開されていることと、一般的なworkstationで高速にresident servingできることは同義ではない。

W32ではこのギャップを極端な形で示すcommunity experimentが大きな反応を得た。Cutoff前の投稿では、expertをdiskから逐次streamするpure-C99 implementationによりpeak RAMを約8.24 GBへ抑えた、という主張が注目された\cite{kimi-k3-x-lowram}。ただしこれはMoonshot公式の推奨要件ではなく、今回独立reproductionも行っていない。

\begin{communitynote}[「動く」と「使える」の間]
Reaction Passでは、約8 GB peak RAMという主張がMain windowで強い注目を集めた一方、Cutoff後には大容量storageと極端に低いthroughputを実用上の制約として指摘する反応も得られた\cite{kimi-k3-x-practicality}。公式repositoryが約1.56 TBであることも踏まえると、この実験の本質は「巨大MoEは少RAMでも実行可能」という勝利宣言ではなく、\textbf{RAM・storage・I/O・latencyをどう交換するか}というsystems trade-offの可視化にある。
\end{communitynote}

\subsection*{W32で見えたのは、modelからecosystemへの評価軸の移動}

3件をまとめて見ると、``open model''という言葉だけでは差が大きすぎる。Qwen3.8-MaxはPreviewへの関心と将来のweights公開期待、DeepSeek-V4-Flash-0731は同一architectureへのagent-oriented post-trainingとMIT公開model、Kimi K3はfrontier規模のfull weights公開と、それを動かす際の物理的な制約をそれぞれ代表している\cite{qwen38-token-plan,deepseek-v4-flash-0731,deepseek-v4-flash-hf,kimi-k3-hf}。

この違いは、frontier modelの評価単位がmodel checkpointだけではなくなったことを示す。\textbf{weights、license、harness、reasoning budget、serving engine、memory/storage hierarchyまで含めて初めて「使えるモデル」になる。}Cutoff後のSGLang v0.5.17がKimi K3とMiniMax H3を同時にserving対象へ取り込んだことも、この方向を補強している\cite{sglang-0517}。次に見るべきなのは「どのモデルが一番強いか」だけでなく、公開された能力がどれだけ速く再現可能なecosystemへ落ちてくるかである。
